% Chapters/3-Quark-Model-and-EFT.tex

\chapter{The Quark Model and Effective Theories of Hadrons}
\label{ch:quark_model_eft}

\section{Motivation: why do we need more than the quark model?}

So far we have used approximate internal symmetries (isospin and flavor SU(3))
to organize hadrons into multiplets. Symmetry, however, does not by itself provide
the \emph{dynamics}: it does not explain \emph{why} some hadrons are much lighter than others,
nor does it systematically predict low-energy interactions.

A first dynamical attempt is the (constituent) non-relativistic quark model:
hadrons are bound states of quasi-particles (constituent quarks) moving slowly,
with masses approximately given by sums of constituent masses plus small interaction corrections.
This picture is qualitatively useful for spectroscopy, but it becomes internally inconsistent
in the light-hadron sector (especially because of the pion).
The goal of this chapter is:
\begin{itemize}
  \item to see \emph{why} the non-relativistic quark model fails for light hadrons,
  \item to understand the pion as a (pseudo-)Goldstone boson of chiral symmetry breaking,
  \item to introduce the EFT logic (chiral perturbation theory) describing low-energy hadron physics.
\end{itemize}

For broader context on the quark model and hadron spectroscopy, see
standard particle-physics texts such as \citep{halzen1984,griffiths2008,thomson2013},
and for QFT background and symmetry concepts see \citep{peskin1995,schwartz2014,weinberg1995qtf1}.

\section{The non-relativistic quark model}

\subsection{Basic assumption and hyperfine interaction}

The simplest dynamical assumption is:
\begin{center}
\emph{Quarks move slowly inside hadrons, so hadron masses are roughly sums of constituent quark masses,
plus relatively small interaction corrections.}
\end{center}

To account for splittings between spin multiplets (e.g.\ $0^-$ vs.\ $1^-$ mesons and
$\tfrac12^+$ vs.\ $\tfrac32^+$ baryons), the effective interaction must depend on spin.
A common schematic ingredient is a hyperfine (spin--spin) Hamiltonian of the form
\begin{equation}
H_{\rm hf} = a \sum_{i\neq j}\frac{\vec S_i\cdot \vec S_j}{m_i m_j},
\label{eq:hf_hamiltonian}
\end{equation}
where $m_i$ are constituent masses, $\vec S_i$ are spin operators,
and $a$ is an effective constant encoding interaction details and wavefunction overlap.
This structure is widely used in introductory discussions of hadron spectroscopy
\citep{halzen1984,griffiths2008,thomson2013}.

\subsection{Mesons: $q_1\bar q_2$}

For mesons, a schematic mass formula is
\begin{equation}
M_{q_1\bar q_2} \;=\; m_1 + m_2
+ \frac{a}{m_1 m_2}\left(\frac{S(S+1)}{2}-\frac{3}{2}\right),
\qquad S=\text{meson spin},
\label{eq:meson_mass_formula}
\end{equation}
where the spin factor comes from $\langle \vec S_1\cdot \vec S_2\rangle$.
In the light SU(2)/SU(3) limit one may take $m_1=m_2=m$, giving
\begin{align}
J^P=0^-:\quad & M = 2m - \frac{3a}{4m^2},
\label{eq:ps_mass}
\\
J^P=1^-:\quad & M' = 2m + \frac{a}{4m^2}.
\label{eq:v_mass}
\end{align}

If we (naively) identify these with $\pi$ and $\rho$, then
\begin{equation}
M'-M \;=\; \frac{a}{m^2} \;\simeq\; m_\rho - m_\pi \;\approx\; 630~\text{MeV}.
\label{eq:meson_split_scale}
\end{equation}
In such a fit the hyperfine splitting scale is not ``small'' compared to typical constituent masses
(one often finds $m$ of order a few hundred MeV).
This already signals a breakdown of the underlying non-relativistic hypothesis in the light sector:
the ``interaction correction'' is comparable to (or larger than) the would-be constituent mass scale.
The core issue is not the specific numbers, but the \emph{self-consistency} of the expansion
assumed by the model \citep{griffiths2008,thomson2013}.

\subsection{Baryons: $q_1 q_2 q_3$}

For baryons a similar schematic expression is
\begin{equation}
M_{q_1 q_2 q_3} \;=\; m_1+m_2+m_3
+a'\sum_{i\neq j}\frac{\langle \vec S_i\cdot \vec S_j\rangle}{m_i m_j}.
\label{eq:baryon_mass_general}
\end{equation}
In the simplified SU(3)-symmetric limit $m_1=m_2=m_3=m$,
the spin dependence can be written as
\begin{equation}
M_{qqq} \;=\; 3m+\frac{a'}{2m^2}\left(S(S+1)-\frac{9}{4}\right),
\qquad S=\text{baryon spin}.
\label{eq:baryon_mass_simplified}
\end{equation}
For $J^P=\tfrac12^+$ (nucleon) and $J^P=\tfrac32^+$ ($\Delta$),
\begin{align}
J^P=\frac12^+:\quad & M = 3m - \frac{3a'}{4m^2},
\\
J^P=\frac32^+:\quad & M' = 3m + \frac{3a'}{4m^2}.
\end{align}
Hence
\begin{equation}
M'-M \;=\; \frac{3a'}{2m^2}.
\label{eq:baryon_split_scale}
\end{equation}
Numerically, the $N$--$\Delta$ splitting is a few hundred MeV, again not parametrically small
relative to typical constituent scales. The model can be a useful organizer of spectra,
but it does not provide a controlled low-energy expansion for light hadrons
\citep{halzen1984,griffiths2008,thomson2013}.

\subsection{Conclusion: inconsistency for light hadrons}

The central message is:
\begin{itemize}
  \item The non-relativistic quark model assumes a small interaction correction to constituent masses.
  \item In the light sector (where the pion is exceptionally light), the fitted ``hyperfine''
  effects are not small compared to the masses being corrected.
\end{itemize}
Therefore, for low-energy QCD we need a different viewpoint. The correct language is:
\emph{symmetries plus effective field theory}. In particular, chiral symmetry and its breaking
control the pion sector and organize interactions as an expansion in momenta and quark masses
\citep{peskin1995,weinberg1995qtf1,schwartz2014,donoghue1992}.

\section{The linear sigma model and chiral symmetry breaking}

\subsection{The pion puzzle}

Experimentally,
\begin{equation}
m_\pi \sim 140~\text{MeV} \ll m_\rho \sim 770~\text{MeV}.
\end{equation}
This suggests that the up and down quark masses are small compared to the QCD scale.
In the limit $m_u\simeq m_d\simeq 0$:
\begin{itemize}
  \item left- and right-handed Dirac fields decouple,
  \item isospin SU(2) enlarges to the \textbf{chiral symmetry}
  \begin{equation}
  SU_L(2)\otimes SU_R(2).
  \end{equation}
\end{itemize}

We do \emph{not} observe approximate parity doublets throughout the hadron spectrum,
so chiral symmetry cannot be realized linearly on the spectrum in an unbroken way.
The standard interpretation is that it is \textbf{spontaneously broken} by the vacuum,
producing (nearly) massless Goldstone bosons: the pions
\citep{peskin1995,weinberg1995qtf1,schwartz2014,donoghue1992}.

\subsection{Mesonic realization}

A convenient mesonic field realizing $SU_L(2)\otimes SU_R(2)$ is a $2\times 2$ matrix
\begin{equation}
M(x)=\sigma(x)I_2+i\,\vec\tau\cdot\vec\phi(x),
\qquad \sigma,\vec\phi\in\mathbb{R},
\label{eq:M_field_def}
\end{equation}
where $\vec\phi$ are pseudoscalars (identified with pions) and $\sigma$ is a scalar.
Consider the Lagrangian
\begin{equation}
\mathcal{L}
=\frac14\,\mathrm{tr}\!\left(\partial_\mu M^\dagger \partial^\mu M\right)
-\frac{m^2}{4}\,\mathrm{tr}\!\left(M^\dagger M\right)
-\lambda\left[\mathrm{tr}\!\left(M^\dagger M\right)\right]^2,
\qquad \lambda>0,
\label{eq:lin_sigma_lagrangian}
\end{equation}
which is invariant under
\begin{equation}
M \to g_L\,M\,g_R^\dagger,
\qquad g_L\in SU_L(2),\ \ g_R\in SU_R(2),
\label{eq:chiral_transform_M}
\end{equation}
with isospin SU(2) corresponding to $g_L=g_R$.
This is the textbook linear sigma model realization of chiral symmetry
\citep{peskin1995,schwartz2014,weinberg1995qtf1}.

\subsection{Vacuum structure and Goldstone bosons}

Define the invariant
\begin{equation}
\rho \equiv \frac12\,\mathrm{tr}(M^\dagger M),
\end{equation}
so for constant fields the potential is
\begin{equation}
V(\rho) = \frac{m^2}{2}\rho + 4\lambda \rho^2.
\end{equation}
If $m^2>0$, the minimum is at $\rho=0$ and scalars/pseudoscalars are degenerate.
If $m^2<0$, the minimum is at
\begin{equation}
0=\frac{m^2}{2}+8\lambda\rho
\qquad\Rightarrow\qquad
\rho_0 = -\frac{m^2}{16\lambda} > 0.
\end{equation}
Isospin invariance selects a vacuum proportional to the identity:
\begin{equation}
M_0=\sqrt{\rho_0}\,I_2 \equiv f_\pi\,I_2.
\label{eq:vev_fpi}
\end{equation}

A convenient parameterization of the vacuum manifold is
\begin{equation}
M(x)=\bigl(f_\pi + S(x)\bigr)\,U(x),
\qquad U(x)\in SU(2),\quad S(x)\in\mathbb{R},
\label{eq:polar_decomp}
\end{equation}
where $U(x)$ contains the Goldstone modes.
If we focus on the Goldstone sector (set $S=0$ at leading order),
the kinetic term becomes
\begin{equation}
\mathcal{L}_{\rm Goldstone}
=\frac{f_\pi^2}{4}\,\mathrm{tr}\!\left(\partial_\mu U^\dagger\partial^\mu U\right),
\label{eq:nonlinear_sigma_kin}
\end{equation}
which is the non-linear sigma model.

Writing
\begin{equation}
U(x)=\exp\!\left(\frac{i\,\vec\pi(x)\cdot\vec\tau}{f_\pi}\right),
\qquad
\vec\pi\cdot\vec\tau=
\begin{pmatrix}
\pi^0 & \sqrt2\,\pi^+ \\
\sqrt2\,\pi^- & -\pi^0
\end{pmatrix},
\qquad
\pi^\pm=\frac{\pi^1\mp i\pi^2}{\sqrt2},
\end{equation}
and expanding shows that the $\pi^a$ are massless when chiral symmetry is exact.
This is the Goldstone pattern
\begin{equation}
SU_L(2)\otimes SU_R(2)\ \longrightarrow\ SU(2)
\qquad\Rightarrow\qquad
3\ \text{Goldstone bosons}.
\end{equation}
For the general QFT mechanism of spontaneous symmetry breaking and Goldstone bosons,
see \citep{peskin1995,schwartz2014,weinberg1995qtf1}.

\subsection{The scalar mode and explicit breaking}

In \eqref{eq:polar_decomp}, the real field $S(x)$ is the heavy scalar mode.
From the potential (for $m^2<0$) one finds a nonzero scalar mass scale, schematically
\begin{equation}
m_S^2 \sim -2m^2 > 0,
\end{equation}
so $S$ can be integrated out at sufficiently low energies.

Small quark masses $m_u\simeq m_d\equiv m_q$ explicitly break chiral symmetry while preserving isospin.
At the level of the effective mesonic theory, one can represent explicit breaking through a term
\begin{equation}
\delta\mathcal{L}
= \frac{f_\pi^2 B}{2}\,\mathrm{tr}\!\left(M_q\,U^\dagger + U\,M_q^\dagger\right),
\qquad M_q=m_q I_2,
\label{eq:explicit_breaking_term}
\end{equation}
where $B$ is a low-energy constant.
Expanding $U=\exp(i\pi^a\tau^a/f_\pi)$ to quadratic order gives the parametric scaling
\begin{equation}
m_\pi^2 \;\propto\; m_q,
\qquad\Rightarrow\qquad
m_\pi \;\propto\; \sqrt{m_q},
\label{eq:mpi_scaling}
\end{equation}
so pions are light because they are pseudo-Goldstone bosons.
This is one of the key low-energy lessons encoded systematically in chiral EFT
\citep{donoghue1992,schwartz2014,thomson2013}.

\section{Including nucleons: why $m_N$ does not vanish as $m_q\to 0$}

Let $\psi=(p,n)^T$ be the nucleon isodoublet, with chiral projections $\psi_{L,R}$.
A direct nucleon mass term $\bar\psi\psi$ violates chiral symmetry.
However, a chiral-invariant Yukawa structure is
\begin{equation}
\mathcal{L}_N
=\bar\psi_L i\slashed{\partial}\psi_L+\bar\psi_R i\slashed{\partial}\psi_R
- d\,\bar\psi_L M \psi_R - d\,\bar\psi_R M^\dagger \psi_L.
\label{eq:nucleon_yukawa}
\end{equation}
When $M$ condenses, $M\to f_\pi I_2$, this generates
\begin{equation}
m_N = d\,f_\pi,
\label{eq:mN_from_fpi}
\end{equation}
which remains finite in the chiral limit. Therefore,
\begin{equation}
m_\pi \to 0 \ \ (m_q\to 0),
\qquad \text{but}\qquad
m_N \sim f_\pi \neq 0.
\end{equation}
Physically, most of the nucleon mass is tied to the scale of spontaneous chiral symmetry breaking,
not directly to the small explicit quark masses
\citep{peskin1995,schwartz2014,donoghue1992}.

At low energies, explicit breaking induces corrections analytic in $m_q$,
so one expects at leading order
\begin{equation}
m_N = d f_\pi + \mathcal{O}(m_q),
\end{equation}
while the pion mass begins already at order $m_q$ as in \eqref{eq:mpi_scaling}.

\section{The non-linear sigma model and chiral perturbation theory (EFT)}

\subsection{Low-energy regime and building blocks}

At momenta well below the hadronic scale, and in particular for processes with
\begin{equation}
p \ll m_N,
\end{equation}
pions can be produced, but most other hadrons cannot. This motivates an EFT built
only out of pion degrees of freedom: \textbf{chiral perturbation theory}.

The building blocks are:
\begin{itemize}
  \item the pion field matrix $U(x)\in SU(2)$, with $U^\dagger U=1$,
  \item an explicit-breaking spurion $M_q$ that transforms as $M_q\to g_L M_q g_R^\dagger$
  and is set at the end to $M_q=m_q I_2$.
\end{itemize}
Power counting is
\begin{equation}
\partial_\mu \sim p \sim m_\pi,
\qquad
M_q \sim m_q \sim m_\pi^2 \sim p^2.
\end{equation}
Then one writes \emph{all} operators consistent with Lorentz symmetry,
chiral symmetry (with spurions), parity, and charge conjugation, organized
by order in $p$. The UV physics is encoded in low-energy constants multiplying
these operators. This is the general EFT philosophy in a concrete hadronic setting
\citep{donoghue1992,schwartz2014,weinberg1995qtf1}.

\subsection{Leading order: only two universal terms}

At leading order, $\mathcal{O}(p^2)$, the pion EFT has the universal form
\begin{equation}
\mathcal{L}^{(2)}_{\chi}
=
\frac{f_\pi^2}{4}\,\mathrm{tr}\!\left(\partial_\mu U^\dagger \partial^\mu U\right)
+
\frac{f_\pi^2 B}{2}\,\mathrm{tr}\!\left(M_q\,U^\dagger + U\,M_q^\dagger\right).
\label{eq:chpt_LO}
\end{equation}
Thus, at very low energies pion physics depends essentially on $f_\pi$ and on one
additional constant $B$ (which sets $m_\pi$ once $M_q$ is fixed).
At next-to-leading order, $\mathcal{O}(p^4)$, many additional operators appear,
each with its own low-energy constant, encoding short-distance QCD effects
\citep{donoghue1992,schwartz2014}.

\section{Pion--nucleon EFT and the axial coupling $g_A$}

\subsection{Chiral-covariant building blocks}

To include nucleons while keeping a controlled expansion for $p\ll m_N$,
it is convenient to introduce $u(x)$ such that
\begin{equation}
U(x)=u(x)^2.
\end{equation}
Define
\begin{equation}
v_\mu \equiv \frac12\left(u\partial_\mu u^\dagger + u^\dagger\partial_\mu u\right),
\qquad
a_\mu \equiv \frac12\left(u\partial_\mu u^\dagger - u^\dagger\partial_\mu u\right),
\label{eq:va_def}
\end{equation}
which transform as chiral-covariant vector/axial combinations.
Then the leading pion--nucleon Lagrangian can be written schematically as
\begin{equation}
\mathcal{L}_{\pi N}^{(1)}
=
\bar N\left(i\slashed{\partial}+\slashed{v}\right)N
+ g_A\,\bar N\,\slashed{a}\gamma^5 N
- m_N\,\bar N N
+ \cdots,
\label{eq:piN_LO}
\end{equation}
where the ellipsis denotes higher-order terms in the derivative/mass expansion.
The coefficient multiplying the axial structure is the \textbf{axial coupling} $g_A$.
It is not fixed by symmetry alone; it must be taken from experiment.
This is a typical EFT feature: symmetry fixes the operator basis and power counting,
but not the numerical values of the low-energy constants
\citep{donoghue1992,schwartz2014,thomson2013}.

\subsection{Very low-energy scattering}

Expanding $u$ in pion fields yields, schematically,
\begin{equation}
a_\mu \simeq \frac{i}{2f_\pi}\,\partial_\mu(\vec\pi\cdot\vec\tau) + \cdots,
\qquad
v_\mu \simeq \mathcal{O}\!\left(\frac{\pi\,\partial\pi}{f_\pi^2}\right),
\end{equation}
so at the lowest energies pion--nucleon interactions are governed primarily by
$f_\pi$ and $g_A$ through the leading terms in \eqref{eq:piN_LO}.
This explains why a small set of measured constants can control a large class of
low-energy processes in the pion sector \citep{donoghue1992,schwartz2014}.

\section{Final remarks and SU(3) generalization}

\begin{itemize}
  \item The heavy scalar mode in the linear sigma model is \emph{model-dependent}.
  The \emph{universal} low-energy content is encoded in the EFT for $U(x)$
  (and in the pion--nucleon EFT if baryons are included).
  \item The general EFT logic (separation of scales, symmetry-based operator expansion,
  and low-energy constants) is part of the standard QFT toolkit
  \citep{peskin1995,schwartz2014,weinberg1995qtf1}.
\end{itemize}

Finally, the non-linear sigma model generalizes to flavor SU(3). In the idealized limit
$m_u\simeq m_d\simeq m_s\simeq 0$ one has
\begin{equation}
SU_L(3)\otimes SU_R(3)\ \longrightarrow\ SU(3),
\end{equation}
with $U(x)\in SU(3)$ written as
\begin{equation}
U(x)=\exp\!\left(\frac{i\,\lambda_a \phi^a}{f}\right),
\end{equation}
so that kaons and the $\eta$ appear as additional approximate Goldstone bosons.
In the isospin limit, explicit breaking can be represented by a mass spurion
\begin{equation}
M_q=
\begin{pmatrix}
m_q & 0 & 0\\
0 & m_q & 0\\
0 & 0 & m_s
\end{pmatrix},
\end{equation}
which generates the observed mass pattern among the pseudo-Goldstone multiplet
\citep{thomson2013,donoghue1992}.
